% !TEX root = ../nc-quadric-classification.tex
% 2026-11-21 / v1.36
\section{明示 noncubic 表示族の相対内在的実現}
\label{sec:relative-noncubic-families}

第~\ref{sec:noncubic-stack}節の中央二重平方族について、
canonical compact 圏と任意の Noetherian 基底変更を含む全相対条件を確認する。
重み付き射影スタック上の Clifford 代数として heart を表し、
その対角加群の perfectness を幾何ファイバーから持ち上げる。
これにより表示スタック $\mathfrak N^{\mathrm{pres}}$ から
内在的 $\mathfrak H_2$ への実際の射を得る。
全ての基底上平坦 coherent 対象が compact であるという追加条件は用いない。
任意の noncubic 特殊ファイバーを持つ相対族の正規形や、
cubic を横断する開近傍の効果性とは区別する。

\subsection{中央有限代数に対する fiberwise perfectness の持上げ}

\begin{lemma}\label{lem:nc-fiber-perfect}
$R\to D$ を Noetherian 可換環の準同型、$H$ を中心に $D$ の作用する
有限 $D$-代数、$M$ を有限 $D$-加群である右 $H$-加群とする。
$H,M$ は $R$ 上平坦と仮定する。
$\mathfrak q\in\operatorname{Spec} D$、$\mathfrak p=\mathfrak q\cap R$ とする。
対応する fiber の点において
\[
 M\otimes_R\kappa(\mathfrak p)
 \quad\text{が}\quad
 H\otimes_R\kappa(\mathfrak p)\text{ 上 perfect}
\]
ならば、$M$ は $\mathfrak q$ のある近傍で $H$ 上 perfect である。
\end{lemma}

\begin{proof}
完全な証明は相対モジュライ論文 \cite[\texttt{relative_noncubic_families}]{QuadricModuli2026} に置く。
\end{proof}

\begin{remark}
この補題は、$R$ 上の平坦性を省いた fiber 判定ではない。
特に有限自由解消の syzygy の $R$-平坦性と、最後の
$\operatorname{Tor}_1^R(K,\kappa(\mathfrak p))=0$ の両方を使う。
基底の被約性、正則性、有限 Krull 次元は用いない。
\end{remark}

\subsection{noncubic heart の weighted stack 表示}

$k$ を標数零の体、$R$ を可換 Noetherian $k$-代数とする。
判別式が可逆な二元二次形式 $f_1,f_2\in R[u,v]_2$ を取り、
\[
 B=R[u,v]\langle c,d\rangle/(c^2-f_1,d^2-f_2),
 \qquad h=cd+dc,
 \qquad C=R[u,v,h],\quad\deg(u,v,h)=(1,1,2)
\]
と置く。$u,v$ は中心にある。
補題~\ref{lem:noncubic-relative-clifford-center} より
\[
 B=C\oplus Cc\oplus Cd\oplus C(cd)
\]
であり、この表示は全ての基底変更と両立する。
$X=\PP_R(1,1,2)=[(\operatorname{Spec} C\setminus V(C_+))/\mathbb G_m]$ とし、
$\mathcal B$ を $B$ の対応する $X$ 上の代数層とする。
\[
 \mathcal B\simeq\OO_X\oplus\OO_X(-1)^{\oplus2}\oplus\OO_X(-2)
 \quad\text{as }\OO_X\text{-modules}.
 \tag{NCF-1}\label{eq:nc-clifford-sheaf}
\]

\begin{proposition}\label{prop:nc-qgr-weighted}
構造対象を保つ完全同値
\[
 \operatorname{QGr} B\simeq\operatorname{QCoh}(X,\mathcal B),\qquad
 \qgr B\simeq\coh(X,\mathcal B)
\]
がある。これは任意の Noetherian 基底変更における module-object 比較と整合する。
coherent 対象間の Hom は有限 $R$-加群であり、任意の平坦基底変更と可換する。
対象・射・完全列について fppf descent が成立する。
\end{proposition}

\begin{proof}
完全な証明は相対モジュライ論文 \cite[\texttt{relative_noncubic_families}]{QuadricModuli2026} に置く。
\end{proof}

\begin{lemma}[weighted cohomology]\label{lem:nc-weighted-cohomology}
任意の $j\in\mathbb Z$ について $\mathbf R\Gamma(X,\OO_X(j))$ は
有限 projective $R$-複体であり、任意の基底変更と可換する。
\[
 H^0(X,\OO_X(j))=C_j,\qquad H^1(X,\OO_X(j))=0,
 \qquad H^2(X,\OO_X(j))=C_{-j-4}^{\vee},
\]
ただし負次数の $C$ は零とする。
従って $\mathbf R\Gamma(X,\mathcal B(j))$ も同じ普遍性を持つ。
特に $n\ge0$ なら
\[
 \mathbf R\Gamma(X,\mathcal B(n))=B_n[0].
 \tag{NCF-2}\label{eq:nc-full-orbit-cohomology}
\]
\end{lemma}

\begin{proof}
完全な証明は相対モジュライ論文 \cite[\texttt{relative_noncubic_families}]{QuadricModuli2026} に置く。
\end{proof}

\subsection{canonical compact 圏と相対 smooth proper 性}

\begin{proposition}\label{prop:nc-compact-generator}
\[
 G=\bigoplus_{i=0}^3\mathcal B(i)
\]
は $D(\operatorname{QCoh}(X,\mathcal B))$ の compact generator である。
その dg 自己準同型代数を $E_R=\mathbf R\End(G)$ とすると、
$E_R$ は $R$ 上 perfect であり、任意の Noetherian $R$-代数 $R'$ に対して
\[
 E_R\otimes_R^{\mathbf L}R'\simeq E_{R'},\qquad
 D(\operatorname{QGr} B)^c\simeq\operatorname{Perf}(E_R)
\]
となる。この同定は実際の module-object 引戻しの導来関手に対して
\[
 D(\operatorname{QGr} B)^c\otimes_{\operatorname{Perf} R}\operatorname{Perf} R'
 \simeq D(\operatorname{QGr} B_{R'})^c
 \tag{NCF-3}\label{eq:nc-perfect-basechange}
\]
を与える。compact 対象は $D^b(\qgr B)$ に属する。
\end{proposition}

\begin{proof}
完全な証明は相対モジュライ論文 \cite[\texttt{relative_noncubic_families}]{QuadricModuli2026} に置く。
\end{proof}

\begin{remark}
$G$ を tilting と主張していない。
例えば負の差に \eqref{eq:nc-clifford-sheaf} を適用すると $H^2$ が現れる。
有限な通常の End 代数に置き換えず、$E_R$ は dg 代数として扱う。
\end{remark}

\begin{proposition}\label{prop:nc-relative-smooth}
上の canonical compact 圏は $R$ 上 smooth proper である。
この結論は任意の Noetherian 基底変更後に成立する。
\end{proposition}

\begin{proof}
完全な証明は相対モジュライ論文 \cite[\texttt{relative_noncubic_families}]{QuadricModuli2026} に置く。
\end{proof}

\subsection{同じ根・全軌道・相対 AS4 被覆}

\begin{proposition}\label{prop:nc-root-ample-orbit}
étale 局所に、$B$ の次数付き自己同型 $\sigma$ で
\[
 \sigma^2=\iota,\qquad
 \iota(u)=u,\quad\iota(v)=v,\quad\iota(c)=-c,\quad\iota(d)=-d
\]
となるものを選べる。
canonical compact 圏の相対 Serre 関手と、heart の完全自己同値は
\[
 S=(-)_\iota(-2)[2],\qquad s=(-)_\sigma(1),\qquad s^2=S^{-1}[2]
\]
を満たす。$(\mathcal B,s)$ は AZ ample であり、
右軌道積の全軌道環は $B^\sigma$ に等しい。
これらは任意の Noetherian 基底変更と整合する。
\end{proposition}

\begin{proof}
完全な証明は相対モジュライ論文 \cite[\texttt{relative_noncubic_families}]{QuadricModuli2026} に置く。
\end{proof}

\begin{proposition}[全相対 AS4 条件]\label{prop:nc-full-relative-as4}
前命題の選択の下で
\[
 A'=R[u,v]\langle c,d\rangle/
 \bigl((c^2-f_1)-\eta(d^2-f_2)\bigr),\qquad
 g'=(c^2-f_1)+\eta(d^2-f_2)
\]
と置く。
$A=(A')^\sigma$ と $g'$ は、定義~\ref{def:relative-h2}の相対 quadratic AS4 被覆の全条件を満たし、
\[
 A/(g')=B^\sigma
\]
である。特に全ての Noetherian 基底変更後も $A$ は左右 Noetherian であり、
左右・双加群 Koszul 解消および正則正規切断は任意の底変更後も完全である。
\end{proposition}

\begin{proof}
完全な証明は相対モジュライ論文 \cite[\texttt{relative_noncubic_families}]{QuadricModuli2026} に置く。
\end{proof}

\subsection{表示スタックからの実際の相対族と残る境界問題}

\begin{theorem}\label{thm:noncubic-relative-effectivity}
第~\ref{sec:noncubic-stack}節の表示スタック
\[
 \mathfrak N^{\mathrm{pres}}=[P/G],\qquad
 P=Q^{\mathrm{nd}}\times Q^{\mathrm{nd}}
\]
上の降下した反 Serre 環 $\mathcal R$ は、任意の有限型 affine test $T$ 上で
\[
 (\mathcal C_T,\mathcal O_T)=(\qgr\mathcal R_T,\pi\mathcal R_T)
\]
という実際の相対 $\mathfrak H_2$ 族を定める。
従って自然な射
\[
 q_{\mathrm{nc}}:\mathfrak N^{\mathrm{pres}}\longrightarrow\mathfrak H_2
\]
が存在し、その全ての幾何ファイバーは noncubic である。
ここでは平方根・AS4 被覆は局所的 witness であり、出力の標識に加えない。
\end{theorem}

\begin{proof}
完全な証明は相対モジュライ論文 \cite[\texttt{relative_noncubic_families}]{QuadricModuli2026} に置く。
\end{proof}

\begin{remark}[noncubic 境界との区別]
$q_{\mathrm{nc}}$ の存在だけでは、noncubic 点を含む開近傍の atlas は得られない。
特に任意の Artin 変形は、閉 fiber が noncubic でも
隣接乗法の四次 minor が非零の冪零元となり得る。
従って二形式表示の次元一のパラメータ族で全ての横断変形を捉えたことにはならない。
任意の相対 $\mathfrak H_2$ 族の noncubic 表示への局所還元、
この射の相対全忠実性・monomorphism 性、および cubic を横断する近傍の効果性は別途証明が必要である。
ここでの対角加群の議論は、明示的な中央有限な Clifford order を持つ族に適用したものであり、
一般の root-family が同じ中央有限表示を持つとは仮定していない。
\end{remark}
